% Full LaTeX bundle: uplift two-batch monitor + federated cross-block FSDS
% Compile from docs/latex:
%   pdflatex fsds_monitoring_full_standalone.tex
%   pdflatex fsds_monitoring_full_standalone.tex  (second pass for references)

\documentclass[11pt]{article}
\usepackage[margin=1in]{geometry}
\usepackage{amsmath,amssymb,booktabs}
\usepackage{hyperref}

\title{FSDS Monitoring --- Full Document\\[0.4em]
\large Part~I: Two-Batch Uplift (LOCO--AUUC, Two-Layer Localization)\quad
Part~II: Communication-Efficient Federated Cross-Block Protocol}
\author{Batch FSDS monitoring and federated feature-block attribution}
\date{}

\begin{document}
\maketitle
\tableofcontents
\newpage

\begin{abstract}
\textbf{Part~I (uplift).}
We formalise uplift model monitoring under a two-batch contract: REF calibrates a frozen ranker $\hat\tau$; LIVE is deployment.
Layer~1 (AUUC, LOCO--AUUC, domain-quintile pairwise comparisons) localises ranking failure; Layer~2 (empirical ATE and mean $\hat\tau$ on LIVE slices) localises effect movement, gated by PO-risk on stacked batches.
Outputs are direct targeting policies (REALLOCATE, REFRESH\_REF, RELEARN) without a hierarchical feature tree.

\textbf{Part~II (federated).}
For vertically partitioned features, clients upload only $O(|F_k|)$ attribution summaries per REF/LIVE window.
The server merges ranks, runs BH-FDR, diagnoses drift type, and emits online feature selection actions.
Five UCI benchmarks show $10^2$--$10^3\times$ uplink efficiency versus naive matrix centralization and localize injected shift in targeted blocks.
\end{abstract}

\part{Two-batch uplift monitoring}

% === BEGIN uplift_fsds_two_batch_formulation.tex ===
% Canonical LaTeX formulation: two-batch (REF vs LIVE) uplift monitoring
% isomorphic to batch FSDS, with AUUC ranking layer + uplift effect layer.
% Compile via uplift_fsds_two_batch_standalone.tex

\section{Two-batch uplift monitoring (FSDS-isomorphic formulation)}

\paragraph{Design rationale (post-hoc localization).}
The post-hoc localization is performed at two levels: one based on AUUC (evaluated on a single batch) and the other based on estimated uplift, followed by pseudo-outcome risk.
This design is highly efficient and beneficial for downstream procedures.
Moreover, the AUC of the RF-domain binary classifier is instrumental for monitoring covariate shift.
These elements collectively produce compound insights that facilitate interpretation.

\paragraph{Monitoring contract.}
Fix two disjoint finite samples representing \textbf{calibration} and \textbf{deployment}:
\[
  \mathcal{D}_{\mathrm{ref}} = \{(X_i,T_i,Y_i)\}_{i\in\mathcal{I}_{\mathrm{ref}}},
  \qquad
  \mathcal{D}_{\mathrm{live}} = \{(X_j,T_j,Y_j)\}_{j\in\mathcal{I}_{\mathrm{live}}}.
\]
Split $\mathcal{D}_{\mathrm{ref}}$ into training and evaluation subsets
$\mathcal{D}_{\mathrm{ref}}^{\mathrm{tr}}\cup\mathcal{D}_{\mathrm{ref}}^{\mathrm{ev}}$ (held out for in-batch AUUC).
The \textbf{uplift ranker} $\hat\tau:\mathbb{R}^p\to\mathbb{R}$ is fit only on $\mathcal{D}_{\mathrm{ref}}^{\mathrm{tr}}$ and then \textbf{frozen} when scoring $\mathcal{D}_{\mathrm{live}}$.
No refit on LIVE is permitted in the monitoring layer except inside LOCO ablations (always retrain on REF train only).

\paragraph{Batch vs treatment.}
Introduce batch indicator $W\in\{0,1\}$ with $W=0$ on $\mathcal{D}_{\mathrm{ref}}$ and $W=1$ on $\mathcal{D}_{\mathrm{live}}$.
This is \textbf{not} the randomised treatment $T\in\{0,1\}$ used in the uplift experiment.
All FSDS-style batch tests use $(X,W)$ or stacked $(X,Y,W)$; uplift curves use $(Y,T,\hat\tau(X))$ within each batch.

\paragraph{Isomorphism to predictive batch FSDS.}
\begin{center}
\small
\begin{tabular}{@{}lll@{}}
\toprule
\textbf{Predictive FSDS} & \textbf{Two-batch uplift monitor} & \textbf{Object} \\
\midrule
Global loss / calibration on LIVE & $\mathrm{AUUC}_{\mathrm{live}}$, $\Delta_{\mathrm{AUUC}}$ & ranking quality \\
Global MMD$^2$ on $X$ & $\widehat{\mathrm{MMD}}^2(X_{\mathrm{ref}},X_{\mathrm{live}})$ & covariate shift \\
Domain classifier AUC$(X)$ & $\widehat{\mathrm{AUC}}_{\mathrm{dom}}^{\mathrm{RF}}(X;W)$ & same \\
Effective sample size (overlap) & $\mathrm{ESS}_{\mathrm{ovlp}}(\hat e(W\mid X))$ & comparable support \\
LOGO-MMD on feature groups & LOGO-MMD$(g)$ on $X$ & localised $X$ shift \\
Leave-one-group-out on outcome model & LOCO--AUUC$(g)$ & localised ranking drivers \\
Permute--refit on batch (PO-risk) & DRPerm on $(X,Y,W)$ & $Y\mid X$ shift \\
Post-hoc population subset & Domain-quintile AUUC on LIVE & user-slice ranking \\
\bottomrule
\end{tabular}
\end{center}

\subsection{Validity gates (batch geometry)}

\paragraph{Covariate shift on $X$.}
Let $\hat{\phi}$ be an RKHS mean embedding distance or unbiased MMD$^2$ estimator:
\[
  \widehat{\mathrm{MMD}}^2
  = \widehat{\mathrm{MMD}}^2\bigl(\{X_i\}_{i\in\mathcal{I}_{\mathrm{ref}}^{\mathrm{ev}}},
  \{X_j\}_{j\in\mathcal{I}_{\mathrm{live}}}\bigr).
\]
Train a domain classifier $\hat d_W:\mathbb{R}^p\to[0,1]$ on stacked labelled pairs $(X,W)$ and report
$\widehat{\mathrm{AUC}}_{\mathrm{dom}}^{\mathrm{RF}}=\widehat{\mathrm{AUC}}(\hat d_W(X),W)$.

\paragraph{Batch propensity overlap.}
On the stacked sample, estimate $\hat e(w\mid x)=\mathbb P(W=1\mid X=x)$ and define overlap weights
$v_i = 1/(\hat e_i(1-\hat e_i))$ with clipping.
The effective sample size fraction is
\[
  \mathrm{ESS}_{\mathrm{ovlp}}
  = \frac{\bigl(\sum_i v_i\bigr)^2}{\sum_i v_i^2}\cdot\frac{1}{n_{\mathrm{ref}}+n_{\mathrm{live}}}.
\]
Monitoring attributions (LOCO--AUUC, slice tables) require $\mathrm{ESS}_{\mathrm{ovlp}}$ above a policy threshold; otherwise emit \textsc{RefreshRef} before subset rules.

\paragraph{Randomisation check on $T$.}
Within each batch $w\in\{0,1\}$, test treatment balance on $T$ (e.g.\ Pearson $\chi^2$ against design ratio $\pi$).
Failure $\Rightarrow$ invalid uplift curves; halt AUUC interpretation.

\subsection{Layer 1: ranking monitor (AUUC)}

\paragraph{Uplift ranker.}
Let $\mathcal{A}$ be a meta-learner factory (T-/X-/R-learner or registry nuisance models).
Fit $\hat\tau$ on $\mathcal{D}_{\mathrm{ref}}^{\mathrm{tr}}$:
\[
  \hat\tau \leftarrow \mathcal{A}\bigl(\mathcal{D}_{\mathrm{ref}}^{\mathrm{tr}}\bigr).
\]
Scores on evaluation windows:
\[
  \hat\tau_i^{\mathrm{ref}} = \hat\tau(X_i),\ i\in\mathcal{I}_{\mathrm{ref}}^{\mathrm{ev}};
  \qquad
  \hat\tau_j^{\mathrm{live}} = \hat\tau(X_j),\ j\in\mathcal{I}_{\mathrm{live}}.
\]

\paragraph{Area under the uplift curve (AUUC).}
Sort users by $\hat\tau$ descending; let $n_t(k)$, $n_c(k)$ be treated/control counts among the top-$k$ fraction.
With cumulative treated/control response sums $S_t(k)$, $S_c(k)$, define the uplift curve
$U(k)=\frac{S_t(k)}{n_t(k)}-\frac{S_c(k)}{n_c(k)}$ (with standard sklift conventions on ties and empty strata).
Then
\[
  \mathrm{AUUC}_{\mathrm{ref}}
  = \int_0^1 U_{\mathrm{ref}}(k)\,dk
  \quad\text{on}\quad \mathcal{D}_{\mathrm{ref}}^{\mathrm{ev}},
\]
and analogously $\mathrm{AUUC}_{\mathrm{live}}$ on $\mathcal{D}_{\mathrm{live}}$.
The \textbf{two-batch ranking alarm} is
\[
  \Delta_{\mathrm{AUUC}} = \mathrm{AUUC}_{\mathrm{ref}} - \mathrm{AUUC}_{\mathrm{live}}.
\]
Large $|\Delta_{\mathrm{AUUC}}|$ under stable gates indicates deployment ranking degradation relative to REF calibration.

\paragraph{LOCO--AUUC (feature-localised ranking).}
Let $\mathcal{G}=\{g_1,\ldots,g_G\}$ be optional feature groups; if $\mathcal{G}=\emptyset$, each single coordinate defines a group.
For each $g\in\mathcal{G}$, let $X^{(-g)}$ denote features with group $g$ removed.
Retrain $\hat\tau^{(-g)} \leftarrow \mathcal{A}(\mathcal{D}_{\mathrm{ref}}^{\mathrm{tr}}, X^{(-g)})$ and recompute AUUC on both windows:
\[
  \mathrm{drop}_{\mathrm{ref}}^{(g)}
  = \mathrm{AUUC}_{\mathrm{ref}}^{\mathrm{full}} - \mathrm{AUUC}_{\mathrm{ref}}^{(-g)},\quad
  \mathrm{drop}_{\mathrm{live}}^{(g)}
  = \mathrm{AUUC}_{\mathrm{live}}^{\mathrm{full}} - \mathrm{AUUC}_{\mathrm{live}}^{(-g)}.
\]
Define $\mathrm{drop\_gap}^{(g)}=\mathrm{drop}_{\mathrm{live}}^{(g)}-\mathrm{drop}_{\mathrm{ref}}^{(g)}$.
Large $|\mathrm{drop}_{\mathrm{live}}^{(g)}|$ localises LIVE ranking sensitivity to group $g$;
the vector $(\mathrm{drop}_{\mathrm{ref}}^{(g)})_g$ vs $(\mathrm{drop}_{\mathrm{live}}^{(g)})_g$ compared by Spearman $\rho$ distinguishes covariate-driven vs concept-driven reallocation of feature importance.

\paragraph{User slices on LIVE (domain quintiles).}
Train batch propensity on stacked covariates:
\[
  \hat s(x) = \mathbb P(W=1\mid X=x)
  \quad\text{fit on}\quad
  \{(X_i,W_i)\}_{i\in\mathcal{I}_{\mathrm{ref}}^{\mathrm{pool}}\cup\mathcal{I}_{\mathrm{live}}},
\]
where $\mathcal{I}_{\mathrm{ref}}^{\mathrm{pool}}$ is a REF pool (eval $+$ optional train cap).
On LIVE only, form quantile bins $Q_1,\ldots,Q_K$ by $\hat s(X_j)$ and compute
\[
  \mathrm{AUUC}_k = \mathrm{AUUC}\bigl(\mathcal{D}_{\mathrm{live}}\cap Q_k\bigr),\qquad k=1,\ldots,K.
\]
For slices $A,B\subseteq\mathcal{I}_{\mathrm{live}}$, pairwise ranking contrast
$\Delta_{\mathrm{AUUC}}^{A,B}=\mathrm{AUUC}_A-\mathrm{AUUC}_B$; policy ranks pairs by $|\Delta_{\mathrm{AUUC}}^{A,B}|$ for cap/continue rules.

\paragraph{Overlap-support slice.}
Let $\mathcal{L}_{\mathrm{ovlp}}=\{j\in\mathcal{I}_{\mathrm{live}}:\ \hat e(W=1\mid X_j)\in[\ell,u]\}$.
Report $\mathrm{AUUC}_{\mathrm{ovlp}}$ on $\mathcal{L}_{\mathrm{ovlp}}$.
If $\mathrm{AUUC}_{\mathrm{ovlp}}\gg \mathrm{AUUC}_{\mathrm{live}}$, global $\Delta_{\mathrm{AUUC}}$ is partially explained by batch mix off the shared support.

\subsection{Layer 2: uplift level on LIVE slices (effect layer)}

\paragraph{Observed slice uplift.}
On each quintile $Q_k\subseteq\mathcal{D}_{\mathrm{live}}$ with sufficient $T$-support,
\[
  \widehat{\tau}^{\mathrm{obs}}_k
  = \bar Y_{T=1,k}-\bar Y_{T=0,k}.
\]

\paragraph{Model slice uplift level.}
Using the same frozen REF ranker,
\[
  \bar{\hat\tau}_k = \frac{1}{|Q_k|}\sum_{j\in Q_k} \hat\tau(X_j).
\]

\paragraph{Pairwise effect contrasts (LIVE only).}
For quintiles $Q_a,Q_b$,
\[
  \Delta\widehat{\tau}^{\mathrm{obs}}_{a,b}
  = \widehat{\tau}^{\mathrm{obs}}_a - \widehat{\tau}^{\mathrm{obs}}_b,
  \qquad
  \Delta\bar{\hat\tau}_{a,b} = \bar{\hat\tau}_a - \bar{\hat\tau}_b.
\]
Layer~2 policy uses $|\Delta\widehat{\tau}^{\mathrm{obs}}_{a,b}|$ to prioritise budget reallocation \emph{conditional on concept rejection} (below).
Alignment of $\bar{\hat\tau}_k$ with $\widehat{\tau}^{\mathrm{obs}}_k$ across $k$ informs \textsc{Relearn} vs \textsc{Reallocate}.

\paragraph{Layer separation (do not conflate).}
Layer~1 measures \emph{sorting quality} of $\hat\tau$ on LIVE; Layer~2 measures \emph{level} of treatment contrast on LIVE slices.
A slice may exhibit high $\widehat{\tau}^{\mathrm{obs}}_k$ and low $\mathrm{AUUC}_k$ simultaneously; caps follow Layer~1, budget shifts on Layer~2 require stronger concept evidence.

\subsection{Concept axis: PO-risk on stacked batches}

\paragraph{Pseudo-outcome risk.}
Stack $\mathcal{D}=\mathcal{D}_{\mathrm{ref}}\cup\mathcal{D}_{\mathrm{live}}$ with batch labels $W$.
Cross-fit nuisances $\hat\mu(X)$, $\hat e(W\mid X)$; form pseudo-outcomes $\hat\eta=(Y-\hat\mu)(W-\hat e)$ and fit $\hat\tau_{\mathrm{po}}(X)$ by regressing $\hat\eta$ on $X$.
The observed PO-risk statistic is $R=\mathbb E[\hat\tau_{\mathrm{po}}(X)^2]$ (sample mean on $\mathcal{D}$).

\paragraph{Permute--refit test (DRPerm).}
For $b=1,\ldots,B$, permute batch labels $W^{(b)}$, refit nuisances and $\hat\tau_{\mathrm{po}}^{(b)}$, store $R^{(b)}$.
Two-sided Monte Carlo $p$-value:
\[
  p = \frac{1+\sum_{b=1}^B \mathbf 1\{ R^{(b)}\ge R\}}{B+1}.
\]
Reject at level $\alpha$ $\Rightarrow$ interpret large $\Delta_{\mathrm{AUUC}}$ with flat $\widehat{\mathrm{MMD}}^2$ as $Y\mid X$ (and treatment interaction) shift; enable Layer~2 REALLOCATE rules on empirical ATE pairs.

\subsection{Regime classification (two axes)}

\begin{center}
\small
\begin{tabular}{@{}p{0.22\linewidth}cc@{}}
\toprule
 & $Y\mid X$ stable ($p>\alpha$) & $Y\mid X$ shifted ($p\le\alpha$) \\
\midrule
$X$ stable (MMD, AUC flat) &
Layer~1: LOCO + quintile AUUC &
Layer~1 + Layer~2 ATE; PO-LOCO optional \\
\midrule
$X$ shifted (MMD or AUC $\uparrow$) &
LOGO-MMD$(g)$, quintile AUUC, overlap &
Stratify REF; gates then LOCO + L2 \\
\bottomrule
\end{tabular}
\end{center}

If $\mathrm{ESS}_{\mathrm{ovlp}}$ is low, suspend LOCO and slice effect rules; emit refresh/match REF.

\subsection{Monitoring map (formal pipeline)}

\begin{enumerate}
  \item \textbf{Gates:} SRM$(T)$ per batch; $\widehat{\mathrm{MMD}}^2$, $\widehat{\mathrm{AUC}}_{\mathrm{dom}}$; $\mathrm{ESS}_{\mathrm{ovlp}}$.
  \item \textbf{Layer~1 global:} $\mathrm{AUUC}_{\mathrm{ref}}$, $\mathrm{AUUC}_{\mathrm{live}}$, $\Delta_{\mathrm{AUUC}}$.
  \item \textbf{Layer~1 local:} LOCO--AUUC$(g)$; $\mathrm{AUUC}_k$; pairwise $|\Delta_{\mathrm{AUUC}}^{A,B}|$; optional $\mathrm{AUUC}_{\mathrm{ovlp}}$.
  \item \textbf{Concept:} DRPerm $p$ on stacked $(X,Y,W)$.
  \item \textbf{Layer~2 local:} $\widehat{\tau}^{\mathrm{obs}}_k$, $\bar{\hat\tau}_k$; pairwise $|\Delta\widehat{\tau}^{\mathrm{obs}}_{a,b}|$.
  \item \textbf{Policy:} map statistics to $\{\textsc{Reallocate},\textsc{HoldFeature},\textsc{RefreshRef},\textsc{Relearn},\textsc{Monitor}\}$.
\end{enumerate}

\paragraph{Policy mapping (examples).}
Let $\tau_{\mathrm{AUUC}}$ be AUUC SLA, $\delta_{\mathrm{mmd}}$ MMD threshold, $\epsilon_{\mathrm{ess}}$ ESS floor.
\begin{align*}
  \mathrm{ESS}_{\mathrm{ovlp}}<\epsilon_{\mathrm{ess}}
  &\Rightarrow \textsc{RefreshRef}.\\
  \mathrm{AUUC}_{\mathrm{live}}<\tau_{\mathrm{AUUC}}
  &\Rightarrow \textsc{Reallocate}\ \text{(global pause/cap)}.\\
  k^\star=\arg\min_k \mathrm{AUUC}_k
  &\Rightarrow \textsc{Reallocate}\ \text{(cap quintile $Q_{k^\star}$)}.\\
  \arg\max_{a,b}|\Delta\widehat{\tau}^{\mathrm{obs}}_{a,b}|,\ p\le\alpha
  &\Rightarrow \textsc{Reallocate}\ \text{(budget toward higher $\widehat{\tau}^{\mathrm{obs}}$ slice)}.\\
  \widehat{\mathrm{MMD}}^2\ge\delta_{\mathrm{mmd}},\ \mathrm{ESS}_{\mathrm{ovlp}}\ \text{OK}
  &\Rightarrow \textsc{HoldFeature}\ \text{(LOGO-MMD / LOCO top group)}.\\
  p\le\alpha,\ \text{systematic } \bar{\hat\tau}_k\not\approx\widehat{\tau}^{\mathrm{obs}}_k
  &\Rightarrow \textsc{Relearn}\ \text{after REALLOCATE trial}.
\end{align*}

\subsection{Two-batch information flow}

\begin{center}
\begin{tabular}{@{}c|c|c@{}}
\toprule
 & REF & LIVE \\
\midrule
Fit $\hat\tau$ & $\mathcal{D}_{\mathrm{ref}}^{\mathrm{tr}}$ & --- \\
AUUC eval & $\mathcal{D}_{\mathrm{ref}}^{\mathrm{ev}}$ & $\mathcal{D}_{\mathrm{live}}$ \\
LOCO retrain & $\mathcal{D}_{\mathrm{ref}}^{\mathrm{tr}}$ (per drop) & score only \\
Train $\hat s,\hat e$ for slices & $\mathcal{I}_{\mathrm{ref}}^{\mathrm{pool}}\cup\mathcal{I}_{\mathrm{live}}$ & slice LIVE rows \\
Empirical ATE & --- & $Q_k\subseteq\mathcal{D}_{\mathrm{live}}$ \\
PO-risk / DRPerm & \multicolumn{2}{c}{stacked $\mathcal{D}_{\mathrm{ref}}\cup\mathcal{D}_{\mathrm{live}}$} \\
\bottomrule
\end{tabular}
\end{center}

\subsection{Cross-batch interpretation and business recommendations}
\label{sec:cross-batch-business}

\paragraph{Paragraph 1 --- What the localized subset means (cross-batch).}
After post-hoc localization you hold \emph{named} subsets: a LOCO feature group $g$ with large $\mathrm{drop}_{\mathrm{live}}^{(g)}$; a LIVE domain quintile $Q_k$ (users with similar $\hat s(x)$); a pairwise ranking contrast $(Q_a,Q_b)$ from Layer~1; optionally an overlap-support band and a Layer~2 pair with large $|\Delta\widehat{\tau}^{\mathrm{obs}}_{a,b}|$.
In the two-batch contract these names are \textbf{not} interchangeable with product segments.
$Q_k$ slices \textbf{LIVE rows only}, while $\hat s$ is trained on REF$\cup$LIVE, so high $k$ marks users who look \textbf{more LIVE-like versus REF}, not necessarily ``high value.''
$\mathrm{AUUC}_k$ judges whether the \textbf{REF-frozen} ranker sorts that slice; $\widehat{\tau}^{\mathrm{obs}}_k$ judges \textbf{observed treat--control lift on today's LIVE traffic} inside the slice---they can diverge (high ATE, low AUUC).
Large $\Delta_{\mathrm{AUUC}}$ means REF holdout ranked better than LIVE under the same $\hat\tau$, not that population uplift collapsed.
Before any subset-driven change, require SRM on $T$; if $\mathrm{ESS}_{\mathrm{ovlp}}$ is low, treat only \textsc{RefreshRef} as valid (extend or match REF, defer caps).
If $\widehat{\mathrm{AUC}}_{\mathrm{dom}}^{\mathrm{RF}}$ or MMD flag covariate shift, read quintiles as \textbf{support/mix} slices; if they are flat and DRPerm rejects, read the same quintiles as \textbf{concept} slices and allow Layer~2 budget moves only under PO-risk.

\paragraph{Paragraph 2 --- How to adjust strategy once the subset is identified.}
Default \textbf{action ladder} on LIVE is \textsc{Reallocate} (cap bad slices or \textbf{narrow support}) $\rightarrow$ re-run the two-batch monitor on the \textbf{next} LIVE window $\rightarrow$ \textsc{Relearn} only if ranking or concept remains broken after that trial.
Full prod $\hat\tau$ retrain is intentionally \emph{last}: most cross-batch alerts are fixed by \textbf{where} you treat, not by immediate refit.
\textbf{(i) Feature subset (LOCO / LOGO-MMD top group $g$).}
Do \textbf{not} drop $g$ from the prod feature pipeline because LOCO reported dependence---that drop is exactly what worsened AUUC in the monitor.
Adjust by \textsc{HoldFeature}: keep $g$ in the model; in \emph{targeting rules}, down-rank or exclude \textbf{segments defined by extreme values of $g$} until REF is refreshed or LOGO-MMD falls; if $g$ drove batch shift, plan REF extension so the shifted support enters calibration rather than hard-suppressing $g$ without data.
\textbf{(ii) User subset --- worst quintile $Q_{k^\star}=\arg\min_k \mathrm{AUUC}_k$.}
\textsc{Reallocate} Layer~1: reduce treat probability or impression cap on $Q_{k^\star}$ (e.g.\ score-based targeting only outside $Q_{k^\star}$); shift volume to the \textbf{better} quintile $Q_b$ from the largest pairwise $|\mathrm{AUUC}_a-\mathrm{AUUC}_b|$ when $Q_b$ meets SLA.
\textbf{(iii) Narrow support} (overlap band $\mathcal{L}_{\mathrm{ovlp}}=\{x:\ \hat e(W\mid x)\in[\ell,u]\}$ when $\mathrm{AUUC}_{\mathrm{ovlp}}\gg \mathrm{AUUC}_{\mathrm{live}}$).
This is a first-class strategy, not a fallback: \textbf{deploy the existing REF $\hat\tau$ only on $\mathcal{L}_{\mathrm{ovlp}}$} (gate scores at scoring or policy time), hold out or uniform-random treat outside the band, and report $\mathrm{AUUC}_{\mathrm{live}}$ \emph{conditional on the narrowed policy} in the next window.
High-$\hat s$ LIVE-like tails that fail AUUC are treated as \textbf{out-of-support} until REF absorbs them over several stable days (\textsc{RefreshRef} in parallel).
Narrow support often restores acceptable ranking \textbf{without} RELEARN when global LIVE AUUC is red but the overlap core is green---typical under covariate shift with intact $Y\mid X$ on the shared support.
\textbf{(iv) Layer~2 pair (high $|\Delta\widehat{\tau}^{\mathrm{obs}}|$) when DRPerm $p\le\alpha$.}
Shift same-day treat/holdout budget toward the higher-empirical-ATE quintile in the top pair \textbf{only if} Layer~1 does not cap that quintile; without PO reject, record the pair as \textsc{Monitor} only.
\textbf{(v) \textsc{Relearn}.}
Schedule prod $\hat\tau$ retrain \textbf{after} at least one REALLOCATE or narrow-support trial if \textbf{(a)} DRPerm still rejects and $\bar{\hat\tau}_k$ vs $\widehat{\tau}^{\mathrm{obs}}_k$ disagree across quintiles (model level wrong under concept shift), or \textbf{(b)} $\mathrm{AUUC}_{\mathrm{ovlp}}$ and capped-quintile AUUC remain below SLA despite narrowed deployment (ranking not recoverable by support alone).
Do \textbf{not} RELEARN when only global $\mathrm{AUUC}_{\mathrm{live}}$ is low but $\mathrm{AUUC}_{\mathrm{ovlp}}$ is healthy---keep narrowed support and widen REF instead.
When $\bar{\hat\tau}_k\approx\widehat{\tau}^{\mathrm{obs}}_k$ but AUUC stays weak, prefer recalibrating $e(T\mid X)$ and segment rules before neural or full-graph relearn.
\textbf{Priority when rules collide:} \textsc{RefreshRef} if ESS fails; then Layer~1 cap; \textbf{narrow support} alongside (not instead of) quintile caps when overlap diagnostic fires; Layer~2 spend last and never on a capped quintile; RELEARN after trial failure, not on first alert.

\paragraph{Narrow support --- how to deploy and when to widen.}
Estimate batch propensity $\hat e(W\mid X)$ on stacked REF$\cup$LIVE (same model as ESS).
Fix a band $[\ell,u]$ (implementation default $[0.15,0.85]$ on LIVE scores $\hat e_j$).
The \textbf{deployable support} is $\mathcal{L}_{\mathrm{ovlp}}=\{j\in\mathcal{I}_{\mathrm{live}}:\ \hat e(W\mid X_j)\in[\ell,u]\}$.
At \textbf{scoring time}, continue to emit $\hat\tau(X)$ for all users if the stack requires it; at \textbf{policy time}, apply uplift targeting only when $j\in\mathcal{L}_{\mathrm{ovlp}}$---users outside receive holdout, control, or a non-uplift policy (document which).
Users with high $\hat s(x)$ in the top domain quintiles but $\hat e(W\mid X)\notin[\ell,u]$ are \textbf{explicitly out-of-support}: do not rank-treat them until REF contains comparable rows (parallel \textsc{RefreshRef}).
\textbf{Trigger to narrow:} $\mathrm{AUUC}_{\mathrm{ovlp}}-\mathrm{AUUC}_{\mathrm{live}}>\delta_{\mathrm{ovlp}}$ (team-tuned, e.g.\ $0.01$--$0.02$) with ESS OK and SRM passed---interpretation: global LIVE ranking is dragged by off-support tails, not necessarily wrong $\hat\tau$ on the core.
\textbf{Next-window success:} $\mathrm{AUUC}_{\mathrm{live}}$ computed under the \emph{actual narrowed policy} meets SLA, or $\mathrm{AUUC}_{\mathrm{ovlp}}$ stays green while global LIVE remains secondary.
\textbf{Widen support gradually:} if domain RF AUC falls and MMD drops over consecutive LIVE windows while $\mathrm{AUUC}_{\mathrm{ovlp}}$ holds, expand $[\ell,u]$ or append prior LIVE into REF and re-estimate $\hat e$; do not widen and RELEARN in the same step---pick support expansion first when $\bar{\hat\tau}_k\approx\widehat{\tau}^{\mathrm{obs}}_k$.

\paragraph{Continue REALLOCATE versus schedule RELEARN --- trial protocol.}
Treat \textsc{Reallocate} and narrow support as a \textbf{controlled trial} on the next LIVE batch(es), not a one-day config tweak.
\textbf{Step A (apply).} Implement the subset actions from Paragraph~2: caps on $Q_{k^\star}$, optional volume shift to $Q_b$, optional $\mathcal{L}_{\mathrm{ovlp}}$ gate, optional Layer~2 budget shift if PO-risk rejected and uncapped.
\textbf{Step B (re-monitor).} Keep REF fixed for the trial; score LIVE with the same frozen $\hat\tau$; recompute gates, $\Delta_{\mathrm{AUUC}}$, $\mathrm{AUUC}_k$, $\mathrm{AUUC}_{\mathrm{ovlp}}$, DRPerm $p$, and quintile $\widehat{\tau}^{\mathrm{obs}}_k$, $\bar{\hat\tau}_k$.
\textbf{Continue REALLOCATE} (extend or tighten, no relearn) if any of the following hold:
(i)~$\mathrm{AUUC}_{\mathrm{live}}$ or policy-conditional AUUC improves versus pre-trial but remains below SLA---tighten cap on $Q_{k^\star}$ or shrink $[\ell,u]$ further;
(ii)~$\mathrm{AUUC}_{\mathrm{ovlp}}$ is above SLA while global LIVE is not---keep narrow support and \textsc{RefreshRef} only;
(iii)~PO-risk $p>\alpha$ and LOCO $\mathrm{drop}_{\mathrm{live}}^{(g)}$ concentrates in one group---continue \textsc{HoldFeature} and segment trims;
(iv)~$\bar{\hat\tau}_k\approx\widehat{\tau}^{\mathrm{obs}}_k$ on capped slices (model tracks data) but ranking is weak---recalibrate $e(T\mid X)$ and hold REALLOCATE; RELEARN is not yet justified.
\textbf{Schedule RELEARN} if after one or two trial LIVE windows \textbf{all} of the following hold:
(1)~SRM and ESS still OK (otherwise \textsc{RefreshRef} first);
(2)~ranking still fails on the \emph{deployed} support: $\mathrm{AUUC}_{\mathrm{ovlp}}<\tau_{\mathrm{AUUC}}$ and/or post-cap $\mathrm{AUUC}_{Q_{k^\star}}<\tau_{\mathrm{AUUC}}$;
(3)~concept evidence persists: DRPerm $p\le\alpha$;
(4)~model--data misalignment: $|\bar{\hat\tau}_k-\widehat{\tau}^{\mathrm{obs}}_k|$ large on multiple quintiles or LOCO Spearman between REF and LIVE drops low---$\hat\tau$ surface or label window is stale, not fixable by caps alone.
\textbf{Do not RELEARN} when only (2) fails but (3)--(4) fail---that is a support/mix problem: narrow further and extend REF.
\textbf{After RELEARN}, promote the new model to REF calibration, slide LIVE into REF history, and reset the action ladder to monitor-only until the next alert.

\paragraph{Implementation.}
\texttt{loco\_auuc.loco\_auuc\_monitor} (Layer~1 LOCO);
\texttt{run\_uplift\_subset\_localization} (gates, quintiles, pairwise, rules);
\texttt{run\_uplift\_subset\_benchmark.py} (four REF/LIVE scenarios).
Empirical subset tables: \verb|docs/latex/uplift_fsds_subset_localization_results.tex|.
Narrative playbook: \verb|docs/FSDS_UPLIFT_TWO_BATCH_LANDING.md|.

% === END uplift_fsds_two_batch_formulation.tex ===


\section{Empirical uplift benchmarks (LOCO--AUUC)}
% === BEGIN uplift_fsds_benchmark_results.tex ===
% Auto-generated from loco_auuc_benchmark.json — 2026-10-01
\paragraph{Empirical monitoring benchmarks.}
Four REF/LIVE scenarios: Hillstrom (real email uplift); Hillstrom with injected covariate shift on $X$;
synthetic covariate shift ($f_0$ offset on LIVE); synthetic concept shift (new driver $f_3$ on LIVE).
Layer~1 reports Online PFI on $X$ (MMD$^2$, domain RF AUC, overlap ESS); Layer~2 is group LOCO--AUUC
(2 blocks, 6 retrains per learner). AUUC matches \texttt{sklift}; $|\Delta\mathrm{sklift}|=0$ on all runs below.
Monitor overlap ESS on stacked REF/LIVE is high ($\gtrsim 0.7$) in most rows---LOCO attributions are gated in.
\subparagraph{hillstrom.}
FSDS on $X$: MMD$^2$=0.0014, $\widehat{\mathrm{AUC}}_{\mathrm{dom}}^{\mathrm{RF}}=$0.9381, ESS$_{\mathrm{ovlp}}=$0.690.
\begin{center}
\small
\begin{tabular}{lrrrrrr}
\toprule Learner & AUUC$_{\mathrm{ref}}$ & AUUC$_{\mathrm{live}}$ & gap & ESS & $\rho$ & LOCO (s) \\
\midrule
registry\_logistic\_x & -0.0004 & 0.0486 & -0.0490 & 0.982 & 1.00 & 0.50 \\
sklearn\_xlearner & -0.0449 & 0.0154 & -0.0603 & 0.982 & 1.00 & 2.77 \\
registry\_rf\_x & -0.0451 & 0.0017 & -0.0468 & 0.982 & 1.00 & 1.82 \\
tarnet & -0.0227 & -0.0131 & -0.0096 & 0.982 & 1.00 & 7.58 \\
dragonnet & -0.0105 & 0.0463 & -0.0568 & 0.982 & 1.00 & 7.76 \\
\bottomrule
\end{tabular}
\end{center}
\subparagraph{hillstrom\_covariate\_drift.}
FSDS on $X$: MMD$^2$=0.0039, $\widehat{\mathrm{AUC}}_{\mathrm{dom}}^{\mathrm{RF}}=$1.0000, ESS$_{\mathrm{ovlp}}=$1.000.
\begin{center}
\small
\begin{tabular}{lrrrrrr}
\toprule Learner & AUUC$_{\mathrm{ref}}$ & AUUC$_{\mathrm{live}}$ & gap & ESS & $\rho$ & LOCO (s) \\
\midrule
registry\_logistic\_x & 0.0152 & -0.0236 & 0.0388 & 1.000 & 1.00 & 0.43 \\
sklearn\_xlearner & 0.0144 & -0.0181 & 0.0325 & 1.000 & 1.00 & 3.02 \\
registry\_rf\_x & -0.0341 & -0.0093 & -0.0248 & 1.000 & 1.00 & 1.75 \\
tarnet & -0.0241 & 0.0284 & -0.0525 & 1.000 & 1.00 & 5.91 \\
dragonnet & 0.0147 & -0.0164 & 0.0311 & 1.000 & 1.00 & 7.02 \\
\bottomrule
\end{tabular}
\end{center}
\subparagraph{synthetic\_covariate.}
FSDS on $X$: MMD$^2$=0.0029, $\widehat{\mathrm{AUC}}_{\mathrm{dom}}^{\mathrm{RF}}=$0.9999, ESS$_{\mathrm{ovlp}}=$0.753.
\begin{center}
\small
\begin{tabular}{lrrrrrr}
\toprule Learner & AUUC$_{\mathrm{ref}}$ & AUUC$_{\mathrm{live}}$ & gap & ESS & $\rho$ & LOCO (s) \\
\midrule
registry\_logistic\_x & -0.0515 & 0.0231 & -0.0746 & 0.715 & 1.00 & 0.64 \\
sklearn\_xlearner & -0.0160 & 0.0112 & -0.0272 & 0.715 & -1.00 & 6.81 \\
registry\_rf\_x & -0.0277 & 0.0386 & -0.0663 & 0.715 & 1.00 & 2.60 \\
tarnet & -0.0392 & 0.0106 & -0.0498 & 0.715 & 1.00 & 4.56 \\
dragonnet & -0.0566 & 0.0297 & -0.0863 & 0.715 & 1.00 & 5.27 \\
\bottomrule
\end{tabular}
\end{center}
\subparagraph{synthetic\_concept.}
FSDS on $X$: MMD$^2$=0.0015, $\widehat{\mathrm{AUC}}_{\mathrm{dom}}^{\mathrm{RF}}=$1.0000, ESS$_{\mathrm{ovlp}}=$0.893.
\begin{center}
\small
\begin{tabular}{lrrrrrr}
\toprule Learner & AUUC$_{\mathrm{ref}}$ & AUUC$_{\mathrm{live}}$ & gap & ESS & $\rho$ & LOCO (s) \\
\midrule
registry\_logistic\_x & 0.0620 & -0.0328 & 0.0949 & 0.909 & 1.00 & 0.59 \\
sklearn\_xlearner & 0.0088 & -0.0026 & 0.0114 & 0.909 & -1.00 & 6.77 \\
registry\_rf\_x & -0.0820 & -0.0080 & -0.0741 & 0.909 & -1.00 & 2.58 \\
tarnet & -0.0590 & 0.0257 & -0.0847 & 0.909 & -1.00 & 4.50 \\
dragonnet & 0.0261 & -0.0554 & 0.0815 & 0.909 & -1.00 & 5.42 \\
\bottomrule
\end{tabular}
\end{center}
\noindent\emph{Takeaway from the tables.}\quad
AUUC$_{\mathrm{live}}$ moves under both injected covariate shift and concept shift (four scenarios above).
Covariate cases: domain RF AUC on $X$ rises toward 1; AUUC on LIVE often drops or changes sign.
Concept case: AUUC gap is large; LOCO rank correlation $\rho$ flips sign across blocks.
Registry meta-learners give similar AUUC to Dragon-Net on Hillstrom with much faster LOCO runs.

% === END uplift_fsds_benchmark_results.tex ===


\section{Subset localization and business rules}
% === BEGIN uplift_fsds_subset_localization_results.tex ===
% Auto-generated from uplift_two_layer_benchmark.json — run_uplift_subset_benchmark.py
\paragraph{Two-batch subset localization (empirical).}
Monitor probe: sklearn X-learner; REF/LIVE splits as in \texttt{benchmark\_datasets}; gates include $\widehat{\mathrm{MMD}}^2$, ESS overlap on $\hat e(W\mid X)$, DRPerm PO-risk ($n_{\mathrm{perm}}=32$--$48$). Layer~1: LOCO--AUUC, domain-quintile AUUC, pairwise ranking. Layer~2: empirical ATE and $\bar{\hat\tau}$ per quintile; L2 REALLOCATE rules require PO-risk reject.
Generated (UTC): 2026-10-01T23:19:37.427430+00:00.

\subparagraph{hillstrom.}
Gates: $\widehat{\mathrm{MMD}}^2=0.0016$, $\mathrm{ESS}_{\mathrm{ovlp}}=0.990$, PO-risk $p=0.7273$ (accept). Diagnosis: \texttt{concept\_drift\_or\_tau\_change}. Regime: \texttt{x\_stable\_gap\_large\_check\_loco\_and\_po\_risk}.
Global Layer~1: $\mathrm{AUUC}_{\mathrm{ref}}=-0.0449$, $\mathrm{AUUC}_{\mathrm{live}}=0.0154$, $\Delta_{\mathrm{AUUC}}=-0.0603$.

\noindent\textbf{LOCO--AUUC (top groups).}
\begin{center}\small
\begin{tabular}{lrrr}
\toprule Group & $\mathrm{drop}_{\mathrm{ref}}$ & $\mathrm{drop}_{\mathrm{live}}$ & $\mathrm{drop\_gap}$ \\
\midrule
block\_0 & -0.0868 & -0.0167 & 0.0701 \\
block\_1 & 0.0179 & -0.0033 & -0.0212 \\
\bottomrule\end{tabular}\end{center}
\noindent Top pairwise AUUC (L1): cap \texttt{domain\_quintile\_3}, continue \texttt{domain\_quintile\_4} ($|\Delta\mathrm{AUUC}|=0.1260$).
\noindent\textbf{LIVE quintiles (L1 AUUC + L2 effect).}
\begin{center}\small
\begin{tabular}{lrrrrr}
\toprule Slice & $n$ & $\mathrm{AUUC}_{\mathrm{live}}$ & $\widehat{\tau}^{\mathrm{obs}}$ & $\bar{\hat\tau}$ & $\hat s$ \\
\midrule
Q1 & 220 & 0.0492 & 0.0613 & 0.0837 & 0.559 \\
Q2 & 260 & 0.0422 & 0.0868 & 0.0856 & 0.701 \\
Q3 & 240 & -0.0557 & 0.1101 & 0.0831 & 0.782 \\
Q4 & 228 & 0.0703 & 0.0789 & 0.0764 & 0.880 \\
Q5 & 252 & -0.0031 & 0.0057 & 0.0638 & 0.961 \\
\bottomrule\end{tabular}\end{center}
\noindent Top pairwise empirical ATE (L2): \texttt{domain\_quintile\_3} vs \texttt{domain\_quintile\_5}, $\Delta\widehat{\tau}^{\mathrm{obs}}=0.1043$ (higher: \texttt{domain\_quintile\_3}).
\noindent\textbf{Business rules emitted.}
\begin{itemize}\itemsep2pt
\item[\textsc{REALLOCATE} (P1).] Cap uplift in domain quintile 'domain\_quintile\_3' (AUUC\_live=-0.0557, n=240); continue in 'domain\_quintile\_4' where ranking remains stronger.
\item[\textsc{REALLOCATE} (P2).] Uplift ranking on LIVE depends on group 'block\_0': do not drop this block from the model without replacement; if trimming segments, keep users where this group's LOCO signal is stable (LOCO AUUC drop\_live=-0.0167).
\item[\textsc{REALLOCATE} (P2).] Uplift ranking on LIVE depends on group 'block\_1': do not drop this block from the model without replacement; if trimming segments, keep users where this group's LOCO signal is stable (LOCO AUUC drop\_live=-0.0033).
\item[\textsc{RELEARN} (P2).] X stable but AUUC gap large: extend labels and schedule prod uplift model relearn after REALLOCATE trial.
\item[\textsc{MONITOR} (P4).] Layer-2 diagnostic (PO-risk not reject): largest slice ATE contrast domain\_quintile\_3 vs domain\_quintile\_5 (ΔATE=0.1043) — do not REALLOCATE on ATE alone; use Layer-1 AUUC rules first.
\end{itemize}

\subparagraph{hillstrom\_covariate\_drift.}
Gates: $\widehat{\mathrm{MMD}}^2=0.0041$, $\mathrm{ESS}_{\mathrm{ovlp}}=0.683$, PO-risk $p=1.0000$ (accept). Diagnosis: \texttt{concept\_drift\_or\_tau\_change}. Regime: \texttt{x\_stable\_gap\_large\_check\_loco\_and\_po\_risk}.
Global Layer~1: $\mathrm{AUUC}_{\mathrm{ref}}=0.0144$, $\mathrm{AUUC}_{\mathrm{live}}=-0.0181$, $\Delta_{\mathrm{AUUC}}=0.0325$.

\noindent\textbf{LOCO--AUUC (top groups).}
\begin{center}\small
\begin{tabular}{lrrr}
\toprule Group & $\mathrm{drop}_{\mathrm{ref}}$ & $\mathrm{drop}_{\mathrm{live}}$ & $\mathrm{drop\_gap}$ \\
\midrule
block\_0 & -0.0012 & -0.0105 & -0.0092 \\
block\_1 & 0.0051 & -0.0065 & -0.0117 \\
\bottomrule\end{tabular}\end{center}
\noindent\textbf{LIVE quintiles (L1 AUUC + L2 effect).}
\begin{center}\small
\begin{tabular}{lrrrrr}
\toprule Slice & $n$ & $\mathrm{AUUC}_{\mathrm{live}}$ & $\widehat{\tau}^{\mathrm{obs}}$ & $\bar{\hat\tau}$ & $\hat s$ \\
\midrule
Q1 & 6 & nan & nan & nan & 0.987 \\
Q2 & 0 & nan & nan & nan & nan \\
Q3 & 0 & nan & nan & nan & nan \\
Q4 & 0 & nan & nan & nan & nan \\
Q5 & 1194 & -0.0184 & 0.0964 & 0.0849 & 1.000 \\
\bottomrule\end{tabular}\end{center}
\noindent\textbf{Business rules emitted.}
\begin{itemize}\itemsep2pt
\item[\textsc{REALLOCATE} (P1).] Pause or cap uplift targeting on LIVE until AUUC\_live (-0.0181) recovers above 0.000.
\item[\textsc{REALLOCATE} (P2).] Uplift ranking on LIVE depends on group 'block\_0': do not drop this block from the model without replacement; if trimming segments, keep users where this group's LOCO signal is stable (LOCO AUUC drop\_live=-0.0105).
\item[\textsc{REALLOCATE} (P2).] Uplift ranking on LIVE depends on group 'block\_1': do not drop this block from the model without replacement; if trimming segments, keep users where this group's LOCO signal is stable (LOCO AUUC drop\_live=-0.0065).
\item[\textsc{RELEARN} (P2).] X stable but AUUC gap large: extend labels and schedule prod uplift model relearn after REALLOCATE trial.
\end{itemize}

\subparagraph{synthetic\_covariate.}
Gates: $\widehat{\mathrm{MMD}}^2=0.0034$, $\mathrm{ESS}_{\mathrm{ovlp}}=0.756$, PO-risk $p=0.3636$ (accept). Diagnosis: \texttt{concept\_drift\_or\_tau\_change}. Regime: \texttt{x\_stable\_gap\_large\_check\_loco\_and\_po\_risk}.
Global Layer~1: $\mathrm{AUUC}_{\mathrm{ref}}=-0.0160$, $\mathrm{AUUC}_{\mathrm{live}}=0.0112$, $\Delta_{\mathrm{AUUC}}=-0.0272$.

\noindent\textbf{LOCO--AUUC (top groups).}
\begin{center}\small
\begin{tabular}{lrrr}
\toprule Group & $\mathrm{drop}_{\mathrm{ref}}$ & $\mathrm{drop}_{\mathrm{live}}$ & $\mathrm{drop\_gap}$ \\
\midrule
block\_1 & -0.0161 & 0.0144 & 0.0305 \\
block\_0 & -0.0034 & 0.0022 & 0.0056 \\
\bottomrule\end{tabular}\end{center}
\noindent Top pairwise AUUC (L1): cap \texttt{domain\_quintile\_2}, continue \texttt{domain\_quintile\_4} ($|\Delta\mathrm{AUUC}|=0.1252$).
\noindent\textbf{LIVE quintiles (L1 AUUC + L2 effect).}
\begin{center}\small
\begin{tabular}{lrrrrr}
\toprule Slice & $n$ & $\mathrm{AUUC}_{\mathrm{live}}$ & $\widehat{\tau}^{\mathrm{obs}}$ & $\bar{\hat\tau}$ & $\hat s$ \\
\midrule
Q1 & 195 & -0.0165 & -0.0212 & 0.0046 & 0.726 \\
Q2 & 200 & -0.0338 & 0.0399 & 0.0089 & 0.815 \\
Q3 & 186 & -0.0163 & 0.0312 & 0.0141 & 0.863 \\
Q4 & 176 & 0.0913 & 0.0487 & 0.0015 & 0.900 \\
Q5 & 243 & 0.0155 & -0.1188 & -0.0114 & 0.948 \\
\bottomrule\end{tabular}\end{center}
\noindent Top pairwise empirical ATE (L2): \texttt{domain\_quintile\_4} vs \texttt{domain\_quintile\_5}, $\Delta\widehat{\tau}^{\mathrm{obs}}=0.1675$ (higher: \texttt{domain\_quintile\_4}).
\noindent\textbf{Business rules emitted.}
\begin{itemize}\itemsep2pt
\item[\textsc{REALLOCATE} (P1).] Cap uplift in domain quintile 'domain\_quintile\_2' (AUUC\_live=-0.0338, n=200); continue in 'domain\_quintile\_4' where ranking remains stronger.
\item[\textsc{REALLOCATE} (P2).] Uplift ranking on LIVE depends on group 'block\_1': do not drop this block from the model without replacement; if trimming segments, keep users where this group's LOCO signal is stable (LOCO AUUC drop\_live=0.0144).
\item[\textsc{RELEARN} (P2).] X stable but AUUC gap large: extend labels and schedule prod uplift model relearn after REALLOCATE trial.
\item[\textsc{REALLOCATE} (P3).] Restrict targeting to overlap-support users on LIVE (AUUC\_live=0.0226 on support band vs global 0.0112).
\item[\textsc{MONITOR} (P4).] Layer-2 diagnostic (PO-risk not reject): largest slice ATE contrast domain\_quintile\_4 vs domain\_quintile\_5 (ΔATE=0.1675) — do not REALLOCATE on ATE alone; use Layer-1 AUUC rules first.
\end{itemize}

\subparagraph{synthetic\_concept.}
Gates: $\widehat{\mathrm{MMD}}^2=0.0020$, $\mathrm{ESS}_{\mathrm{ovlp}}=0.991$, PO-risk $p=0.0303$ (reject). Diagnosis: \texttt{concept\_drift\_or\_tau\_change}. Regime: \texttt{x\_stable\_gap\_large\_check\_loco\_and\_po\_risk}.
Global Layer~1: $\mathrm{AUUC}_{\mathrm{ref}}=0.0088$, $\mathrm{AUUC}_{\mathrm{live}}=-0.0026$, $\Delta_{\mathrm{AUUC}}=0.0114$.

\noindent\textbf{LOCO--AUUC (top groups).}
\begin{center}\small
\begin{tabular}{lrrr}
\toprule Group & $\mathrm{drop}_{\mathrm{ref}}$ & $\mathrm{drop}_{\mathrm{live}}$ & $\mathrm{drop\_gap}$ \\
\midrule
block\_1 & -0.0216 & 0.0704 & 0.0921 \\
block\_0 & 0.0024 & -0.0297 & -0.0321 \\
\bottomrule\end{tabular}\end{center}
\noindent Top pairwise AUUC (L1): cap \texttt{domain\_quintile\_3}, continue \texttt{domain\_quintile\_2} ($|\Delta\mathrm{AUUC}|=0.0922$).
\noindent\textbf{LIVE quintiles (L1 AUUC + L2 effect).}
\begin{center}\small
\begin{tabular}{lrrrrr}
\toprule Slice & $n$ & $\mathrm{AUUC}_{\mathrm{live}}$ & $\widehat{\tau}^{\mathrm{obs}}$ & $\bar{\hat\tau}$ & $\hat s$ \\
\midrule
Q1 & 171 & -0.0194 & 0.0389 & 0.0632 & 0.737 \\
Q2 & 143 & 0.0514 & 0.0727 & 0.0891 & 0.782 \\
Q3 & 200 & -0.0408 & -0.1114 & 0.0607 & 0.806 \\
Q4 & 261 & -0.0006 & -0.0802 & 0.0565 & 0.836 \\
Q5 & 225 & 0.0186 & 0.0421 & 0.0537 & 0.882 \\
\bottomrule\end{tabular}\end{center}
\noindent Top pairwise empirical ATE (L2): \texttt{domain\_quintile\_2} vs \texttt{domain\_quintile\_3}, $\Delta\widehat{\tau}^{\mathrm{obs}}=0.1841$ (higher: \texttt{domain\_quintile\_2}).
\noindent\textbf{Business rules emitted.}
\begin{itemize}\itemsep2pt
\item[\textsc{REALLOCATE} (P1).] Pause or cap uplift targeting on LIVE until AUUC\_live (-0.0026) recovers above 0.000.
\item[\textsc{REALLOCATE} (P1).] Cap uplift in domain quintile 'domain\_quintile\_3' (AUUC\_live=-0.0408, n=200); continue in 'domain\_quintile\_2' where ranking remains stronger.
\item[\textsc{REALLOCATE} (P1).] Concept drift (PO-risk): observed uplift differs by slice — higher empirical ATE in 'domain\_quintile\_2' vs pairwise partner (ΔATE=0.1841 on empirical\_ate); shift today's targeting budget toward the high-uplift slice afte\ldots
\item[\textsc{REALLOCATE} (P2).] Uplift ranking on LIVE depends on group 'block\_1': do not drop this block from the model without replacement; if trimming segments, keep users where this group's LOCO signal is stable (LOCO AUUC drop\_live=0.0704).
\item[\textsc{REALLOCATE} (P2).] Uplift ranking on LIVE depends on group 'block\_0': do not drop this block from the model without replacement; if trimming segments, keep users where this group's LOCO signal is stable (LOCO AUUC drop\_live=-0.0297).
\item[\textsc{RELEARN} (P2).] X stable but AUUC gap large: extend labels and schedule prod uplift model relearn after REALLOCATE trial.
\end{itemize}

% === END uplift_fsds_subset_localization_results.tex ===


\section{Operational subset insights}
% === BEGIN uplift_fsds_subset_insights.tex ===
% Comprehensive subset-insight logic for uplift monitoring (tabular FSDS use case).
% No hierarchical feature tree required. \input from uplift_fsds_loco_auuc.tex or standalone.
% Requires: amsmath, booktabs

\section{Subset insights for uplift monitoring}

\paragraph{Two axes (everything else follows from these).}
Monitoring splits into \textbf{(A) did $P(X)$ or the REF/LIVE mixture on $X$ change?} and \textbf{(B) did the outcome mechanism $Y\mid X$ (and treatment interaction) change?}
They drive different signals, different trustworthy subset tools, and different business rules.
AUUC is sensitive to \emph{both}, but \emph{cannot} tell them apart alone---FSDS on $X$ answers (A); PO-risk / DRPerm on stacked batches answers (B).
When (A) and (B) co-occur, run validity gates first (ESS), then localize features (LOCO, LOGO-MMD) and slices (quintiles) before RELEARN.

\begin{center}
\small
\begin{tabular}{@{}p{0.18\linewidth}p{0.36\linewidth}p{0.36\linewidth}@{}}
\toprule
 & \textbf{$Y\mid X$ stable} & \textbf{$Y\mid X$ shifted (PO-risk reject)} \\
\midrule
\textbf{$X$ stable}\newline (MMD/AUC flat) &
Ranking/$\hat\tau$ issue on same support:\newline LOCO--AUUC + quintiles trusted;\newline gap $\approx$ pure targeting drift &
Concept + ranking:\newline LOCO names feature drivers;\newline PO-LOCO names outcome drivers;\newline RELEARN after REALLOCATE \\
\midrule
\textbf{$X$ shifted}\newline (MMD/AUC $\uparrow$) &
Covariate shift:\newline LOGO-MMD (groups), quintile AUUC;\newline AUUC may flip on LIVE slices;\newline recalibrate $e(T\mid X)$, cap segments &
Mixed shift:\newline stratify REF; subset tools on matched support;\newline do not trust global LOCO until ESS OK \\
\bottomrule
\end{tabular}
\end{center}

\paragraph{Goal.}
Global $\mathrm{AUUC}_{\mathrm{live}}$ or $\Delta_{\mathrm{AUUC}}$ flags a problem; \textbf{subset localization} answers \emph{where} on LIVE (which features or which users/slices).
This mirrors distribution-shift FSDS for predictive models, with \textbf{MSE replaced by AUUC} on the ranking layer.
We do \textbf{not} use a hierarchical feature tree unless a known path is required---post-hoc localization uses flat LOCO, quintiles, and pairwise comparisons only.

\subsection{Two monitoring layers: AUUC vs uplift (when to use which)}

Monitoring is intentionally \textbf{two-layered}. Both layers can use the same LIVE slices (domain quintiles $Q_1,\ldots,Q_5$), but they measure different objects.
Do not treat a drop in AUUC as ``uplift went down'' without Layer~2; do not treat a spike in empirical ATE as ``ranking is fine'' without Layer~1.

\paragraph{Layer 1 --- AUUC movement (ranking / targeting).}
\textbf{Fixed object:} REF-fitted scores $\hat\tau(x)$ evaluated on LIVE.
\begin{itemize}
  \item \textbf{Global:} $\mathrm{AUUC}_{\mathrm{live}}$, $\Delta_{\mathrm{AUUC}}=\mathrm{AUUC}_{\mathrm{ref}}-\mathrm{AUUC}_{\mathrm{live}}$ --- SLA alarm: is today's traffic \emph{sorted} well by yesterday's ranker?
  \item \textbf{Features:} LOCO--AUUC $\mathrm{drop}_{\mathrm{live}}^{(g)}$ --- which columns does LIVE ranking \emph{depend on}?
  \item \textbf{Users:} quintile $\mathrm{AUUC}_k$, pairwise $|\mathrm{AUUC}_A-\mathrm{AUUC}_B|$ --- where does the \emph{same} $\hat\tau$ rank badly?
\end{itemize}
AUUC moves under covariate shift, concept drift, and stale $\hat\tau$; FSDS on $X$ and PO-risk disambiguate \emph{why} the alarm fired.
\textbf{Use Layer~1} for scheduled monitor SLAs, whenever $\Delta_{\mathrm{AUUC}}$ fails, and whenever the ops question is \emph{cap/continue targeting by rank} (REALLOCATE on low-AUUC slices, overlap-only scores).

\paragraph{Layer 2 --- Uplift movement (effect level).}
\textbf{Object:} how large treatment benefit is on each slice---not whether $\hat\tau$ orders users correctly.
\begin{itemize}
  \item \textbf{Observed:} $\widehat{\tau}^{\mathrm{obs}}_k=\bar Y_{T=1,k}-\bar Y_{T=0,k}$ on LIVE quintile $k$ (empirical ATE).
  \item \textbf{Model level:} $\bar{\hat\tau}_k=\mathbb E[\hat\tau(X)\mid Q_k]$ on LIVE (mean estimated uplift).
  \item \textbf{Pairwise effect:} $|\widehat{\tau}^{\mathrm{obs}}_a-\widehat{\tau}^{\mathrm{obs}}_b|$ and optionally $|\bar{\hat\tau}_a-\bar{\hat\tau}_b|$.
\end{itemize}
\textbf{Use Layer~2} when PO-risk rejects (concept axis), when ops asks ``where did \emph{realized} uplift move today?'', or to compare $\bar{\hat\tau}_k$ vs $\widehat{\tau}^{\mathrm{obs}}_k$: agreement $\Rightarrow$ reallocate budget toward high-observed slices; systematic mismatch $\Rightarrow$ RELEARN.
Layer~2 is \textbf{not} a substitute for Layer~1 SLA: small slices make ATE noisy; ranking can fail while average effects look flat.

\begin{center}
\small
\begin{tabular}{@{}p{0.28\linewidth}ccp{0.32\linewidth}@{}}
\toprule
\textbf{Situation} & \textbf{L1 AUUC} & \textbf{L2 uplift} & \textbf{Primary read} \\
\midrule
Routine, all green & global & optional spot-check & monitor \\
$\Delta_{\mathrm{AUUC}}\uparrow$, MMD$\uparrow$, ESS OK & LOCO + pairwise AUUC & optional ATE (support confounds) & ranking on shifted support \\
$\Delta_{\mathrm{AUUC}}\uparrow$, MMD flat, ESS OK & LOCO + pairwise AUUC & ATE + $\bar{\hat\tau}$; PO-risk & ranking and/or concept \\
PO-risk reject, MMD flat & yes (stale ranker) & pairwise $\widehat{\tau}^{\mathrm{obs}}$, $\bar{\hat\tau}$ & effect moved; REALLOCATE \\
AUUC OK, PO-risk reject & watch LOCO & Layer~2 drives ops & silent concept \\
ESS low & defer LOCO & defer slice ATE & REFRESH\_REF \\
Overlap AUUC $\gg$ global & band + pairwise AUUC & on overlap only & narrow targeting \\
\bottomrule
\end{tabular}
\end{center}

\paragraph{Combined workflow.}
Gates (SRM, ESS, MMD) $\rightarrow$ Layer~1 global + LOCO + quintile/pairwise AUUC $\rightarrow$ regime (X vs Y$|$X) $\rightarrow$ Layer~2 slice uplift + pairwise ATE (and $\bar{\hat\tau}$) when concept or budget question $\rightarrow$ actions: rank-based REALLOCATE from L1, effect-based REALLOCATE from L2, RELEARN when PO-risk persists and $\bar{\hat\tau}\not\approx\widehat{\tau}^{\mathrm{obs}}$ across slices.
Markdown reference: \verb|docs/FSDS_UPLIFT_TWO_LAYERS.md|.

\paragraph{Two-batch landing (REF vs LIVE).}
All monitoring assumes \textbf{two batches}: REF calibrates $\hat\tau$; LIVE is scored with REF-frozen weights.
Batch indicator $W\in\{0,1\}$ (REF vs LIVE) is \textbf{not} treatment $T$.
FSDS gates (MMD, domain AUC, ESS on $\hat e(W|X)$), LOCO--AUUC (retrain only on REF), domain quintiles (train $\mathbb P(W{=}1\mid X)$ on REF$\cup$LIVE, slice \textbf{LIVE}), and PO-risk (stacked REF$\cup$LIVE) share this contract.
Deployment playbook: \verb|docs/FSDS_UPLIFT_TWO_BATCH_LANDING.md|.

\paragraph{Windows and notation.}
REF (calibration) vs LIVE (deployment); batch $W\in\{0,1\}$; uplift treatment $T$; meta-learner $\hat\tau$ fit on REF train only.
Optional feature \textbf{groups} $\mathcal{G}=\{g_1,\ldots,g_G\}$; if $\mathcal{G}=\emptyset$, LOCO runs \textbf{per feature} (one column dropped at a time).

\subsection{Workflow (order matters when mixture on $X$ is stable)}

\begin{enumerate}
  \item \textbf{Validity.} SRM on $T$; FSDS globals on $X$: $\widehat{\mathrm{MMD}}^2$, $\widehat{\mathrm{AUC}}_{\mathrm{dom}}^{\mathrm{RF}}(X)$; ESS$_{\mathrm{ovlp}}(\hat e(W|X))$ before trusting attributions.
  \item \textbf{Global ranking.} $\mathrm{AUUC}_{\mathrm{ref}}$, $\mathrm{AUUC}_{\mathrm{live}}$, $\Delta_{\mathrm{AUUC}}=\mathrm{AUUC}_{\mathrm{ref}}-\mathrm{AUUC}_{\mathrm{live}}$.
  \item \textbf{LOCO--AUUC} (always when $\Delta_{\mathrm{AUUC}}$ or $\mathrm{AUUC}_{\mathrm{live}}$ fails SLA):
  for each dropped set $g$ (group or single feature $j$), retrain $\hat\tau$ on REF, recompute AUUC on LIVE:
  \[
    \mathrm{drop}_{\mathrm{live}}^{(g)} = \mathrm{AUUC}_{\mathrm{live}}^{\mathrm{full}} - \mathrm{AUUC}_{\mathrm{live}}^{(-g)}.
  \]
  Large $|\mathrm{drop}_{\mathrm{live}}^{(g)}|$ $\Rightarrow$ ranking on LIVE \emph{depends} on that feature subset.
  \item \textbf{Subset localization on LIVE} (after Step 3 when global alert fires):
  \begin{itemize}
    \item \emph{Population slices:} domain-quintile AUUC---partition LIVE by $\hat s(x)=\mathbb P(W{=}1\mid X=x)$; AUUC on each quintile $Q_1,\ldots,Q_5$.
    \item \emph{Pairwise AUUC:} for slices $A,B$, report $\Delta_{A,B}=\mathrm{AUUC}_A-\mathrm{AUUC}_B$; rank pairs by $|\Delta_{A,B}|$ for business rules (cap the worse slice).
    \item \emph{Support slice (optional):} AUUC on LIVE with $\hat e(W|X)\in[\ell,u]$ (overlap band) vs global LIVE.
  \end{itemize}
  \item \textbf{Optional $Y\mid X$ (concept).} PO-risk / DRPerm on stacked REF$\cup$LIVE when Step~1 shows $X$ stable but Step~2 gap is large.
  \item \textbf{Business rules.} Plain sentences: REALLOCATE (cap quintile / overlap-only targeting), HOLD\_FEATURE (high LOCO or LOGO-MMD group), REFRESH\_REF (ESS low), RELEARN (concept after REALLOCATE trial).
\end{enumerate}

\paragraph{When covariate mixture on $X$ is stable ($\widehat{\mathrm{MMD}}^2$ and domain AUC flat).}
User mixture on inputs looks unchanged; a large $\Delta_{\mathrm{AUUC}}$ is then read as \textbf{ranking or concept drift}, not a new population.
\textbf{Do LOCO--AUUC first} (Step 3) to localize \emph{feature}-level drivers; \textbf{then} quintile + pairwise subset localization (Step 4) to localize \emph{who} on LIVE still ranks poorly under the \emph{same} $\hat\tau$.
Add PO-risk (Step 5) if LOCO points to outcome-side change.
Quintiles remain useful because $\hat\tau$ ranking can differ by slice even when $P(X)$ is stable.

\paragraph{When covariate mixture shifts ($\widehat{\mathrm{MMD}}^2$ or domain AUC $\uparrow$).}
Run LOGO-MMD on groups (if $\mathcal{G}\neq\emptyset$): drop group $g$, recompute MMD on $X\setminus g$; large $\mathrm{MMD}_{\mathrm{full}}-\mathrm{MMD}_{\setminus g}$ $\Rightarrow$ group $g$ carries shift on $X$.
Then LOCO--AUUC and quintile AUUC as above (ranking may break on shifted support).

\subsection{Two complementary subset errors (ranking layer)}

\begin{center}
\begin{tabular}{@{}p{0.22\linewidth}p{0.34\linewidth}p{0.36\linewidth}@{}}
\toprule
\textbf{Insight type} & \textbf{Mechanism} & \textbf{Ops read} \\
\midrule
Feature / LOCO--AUUC &
Per-feature or per-group drop, retrain on REF &
Which columns \emph{not} to drop from targeting; which block to refresh \\
Population / quintile AUUC &
AUUC on LIVE slice $Q_k$ &
Which user segment to cap or continue \\
Pairwise AUUC &
$\arg\max_{A,B}|\mathrm{AUUC}_A-\mathrm{AUUC}_B|$ &
Direct A vs B rule (e.g.\ cap $Q_4$, keep $Q_5$) \\
Overlap-band AUUC &
AUUC on comparable-support LIVE &
Global red but core support green $\Rightarrow$ narrow targeting \\
\bottomrule
\end{tabular}
\end{center}

\subsection{Regime summary (what to run first)}

\begin{center}
\small
\begin{tabular}{@{}lll@{}}
\toprule
Signal & Regime & Primary subset tools \\
\midrule
ESS$_{\mathrm{ovlp}}$ low & Mix / incomparable & Refresh REF; no LOCO-driven cuts \\
MMD or domain AUC $\uparrow$, gap $\uparrow$ & Covariate + ranking & LOGO-MMD (if groups), quintile + pairwise \\
MMD flat, gap $\uparrow$ & Mixture stable, ranking/concept & \textbf{LOCO--AUUC first}, then quintile + pairwise; PO-risk \\
\bottomrule
\end{tabular}
\end{center}

\subsection{Concept drift (PO-risk): where uplift moved today}

\paragraph{When to enter this branch.}
Run this \emph{after} validity (SRM, ESS$_{\mathrm{ovlp}}$) and \emph{after} you read Axis~A as flat: $\widehat{\mathrm{MMD}}^2$ and domain AUC on $X$ are not elevated, yet $\Delta_{\mathrm{AUUC}}$ fails SLA.
Then PO-risk / DRPerm on stacked REF$\cup$LIVE (batch indicator $W$) \textbf{rejects} the null that $P(Y\mid X)$ is unchanged across windows.
That is the concept axis: the \textbf{observed treatment--control contrast} on LIVE can move even when the covariate mixture looks stable.

\paragraph{Two pairwise layers (do not conflate).}
\begin{itemize}
  \item \textbf{Pairwise AUUC} (quintiles $Q_a,Q_b$): compares \emph{ranking quality} of the fixed REF meta-learner $\hat\tau$ on each slice.
  Large $|\mathrm{AUUC}_{Q_a}-\mathrm{AUUC}_{Q_b}|$ answers ``where does targeting \emph{sort} users badly today?''---same as the LOCO--AUUC story, but on \emph{users} not features.
  \item \textbf{Pairwise empirical ATE} (same quintiles): on LIVE only, per slice
  $\widehat{\tau}^{\mathrm{obs}}_k = \bar Y_{T=1,k}-\bar Y_{T=0,k}$.
  Pairwise $\widehat{\tau}^{\mathrm{obs}}_a-\widehat{\tau}^{\mathrm{obs}}_b$ answers ``\textbf{where did observed uplift increase or drop today?}''---a direct post-hoc localization of \emph{effect movement}, not model ranking.
\end{itemize}
Optional cross-check: pairwise on $\mathbb E[\hat\tau(X)\mid Q_k]$ flags slices where the \emph{model} thinks uplift moved vs slices where \emph{data} moved (divergence $\Rightarrow$ relearn; alignment $\Rightarrow$ reallocate budget).

\paragraph{Strategy (ordered).}
\begin{enumerate}
  \item \textbf{Confirm concept.} PO-risk reject + flat $X$-FSDS; if ESS low, stop---refresh REF first.
  \item \textbf{LOCO--AUUC} on LIVE: which features still drive $\hat\tau$ ranking (targeting-side drivers).
  \item \textbf{PO-LOCO} (optional): which features drive $Y\mid X$ shift vs targeting LOCO (outcome-side drivers).
  \item \textbf{Slice uplift on LIVE:} partition users by $\hat s(x)=\mathbb P(W{=}1\mid X)$ quintiles; report $\mathrm{AUUC}_k$, $\widehat{\tau}^{\mathrm{obs}}_k$, $\bar{\hat\tau}_k$ per $Q_k$.
  \item \textbf{Pairwise empirical ATE:} sort all $(Q_a,Q_b)$ by $|\widehat{\tau}^{\mathrm{obs}}_a-\widehat{\tau}^{\mathrm{obs}}_b|$; the top pair names the strongest ``who gained vs who lost'' contrast \emph{on today's LIVE traffic}.
  \item \textbf{Act:} REALLOCATE---shift same-day holdout or promo intensity toward the higher-$\widehat{\tau}^{\mathrm{obs}}$ slice \emph{after} SRM and overlap checks; cap the low-uplift partner from the top pairwise row.
  Schedule \textbf{RELEARN} only if PO-risk persists and $\bar{\hat\tau}_k$ vs $\widehat{\tau}^{\mathrm{obs}}_k$ disagree across slices (model no longer tracks realized effects).
\end{enumerate}

\paragraph{Ops read (one sentence).}
Under concept drift, global AUUC is the alarm; \textbf{pairwise empirical ATE on domain quintiles} is the post-hoc map of \emph{which user subset's uplift jumped today}; pairwise AUUC says where the \emph{existing ranker} still fails on those users.

\subsection{Cross-batch business recommendations (interpretation)}

Once a subset is localized (LOCO group, quintile $Q_k$, pairwise contrast, overlap band, or ATE pair), strategy adjustment is spelled out in \textbf{two paragraphs} in Section~\ref{sec:cross-batch-business}: (1) what the subset means across REF/LIVE; (2) concrete \textsc{Reallocate}/\textsc{HoldFeature}/\textsc{RefreshRef}/\textsc{Relearn} moves per subset type and priority when rules collide.

\subsection{Main insights (paper-ready)}

\begin{itemize}
  \item $\mathrm{AUUC}_{\mathrm{live}}$ and $\Delta_{\mathrm{AUUC}}$ are \textbf{sensitive to both covariate shift and concept drift}; FSDS on $X$ and PO-risk separate the regime.
  \item \textbf{LOCO--AUUC} localizes error to \emph{features} (groups optional; default per-column if no groups).
  \item \textbf{Domain-quintile AUUC} and \textbf{pairwise comparisons} localize error to \emph{population slices} on LIVE.
  \item When \textbf{$X$-mixture is stable}, interpret large gap via LOCO--AUUC $\rightarrow$ subset quintiles $\rightarrow$ optional PO-risk---not via hierarchical trees.
  \item Output is \textbf{business rules} (REALLOCATE / HOLD\_FEATURE / REFRESH\_REF / RELEARN), not a second model stack.
\end{itemize}

\noindent Implementation: \texttt{loco\_auuc.loco\_auuc\_monitor} (LOCO); \texttt{run\_uplift\_subset\_localization} (quintiles, pairwise, rules).
Run \verb|Python/run_uplift_subset_benchmark.py| for four-scenario business narratives:
\verb|artifacts/uplift_two_layer_business_insights.md|.
Empirical AUUC tables: \verb|docs/latex/uplift_fsds_benchmark_results.tex|.

% === END uplift_fsds_subset_insights.tex ===


\part{Communication-efficient federated FSDS}

% === BEGIN federated_fsds_comprehensive_preamble.tex ===
% Static methodology — included by federated_fsds_comprehensive_standalone.tex
% Empirical tables: % === BEGIN federated_fsds_protocol_results.tex ===
% Auto-generated — demo_federated_protocol_datasets.py
\section{Empirical validation across tabular benchmarks}

\paragraph{Summary (communication + localization).}
Each run simulates vertical FL: REF/LIVE split, covariate injection on a contiguous feature block, then the three-phase protocol with BH-FDR ($\alpha{=}0.05$). \textbf{Eff} is the ratio of raw centralized matrix elements $(n_{\mathrm{ref}}+n_{\mathrm{live}})\times p$ to total uplink scalars $\sum_k(3|F_k|+3)$. \textbf{Top-1 hit} indicates whether the highest global attribution score falls inside the injected shift block.

\begin{center}
\small
\begin{tabular}{@{}lrrrrrrl@{}}
\toprule
Dataset & $p$ & nodes & FDR sig & Eff & uplink KiB & top-1 & hit \\
\midrule
WDBC\_breast\_cancer & 30 & 5 & 1 & 163$\times$ & 0.82 & \texttt{f10} & yes \\
UCI\_wine & 13 & 3 & 1 & 48$\times$ & 0.38 & \texttt{f7} & yes \\
UCI\_diabetes & 10 & 3 & 1 & 113$\times$ & 0.30 & \texttt{f6} & yes \\
UCI\_ionosphere & 34 & 5 & 1 & 102$\times$ & 0.91 & \texttt{f12} & yes \\
UCI\_heart\_statlog & 13 & 3 & 2 & 73$\times$ & 0.38 & \texttt{f8} & yes \\
\bottomrule
\end{tabular}
\end{center}

\subparagraph{WDBC\_breast\_cancer.}
$n_{\mathrm{ref}}+n_{\mathrm{live}}=341+228$, shift columns $[10,\,16)$. Communication: 105 floats uplink (0.82 KiB) vs 17070 raw matrix floats (0.130 MiB if centralized)---\textbf{163$\times$} fewer scalars over the wire. FDR significant: 1; actions: \textsc{keep}=29, \textsc{candidate\\_retire}=1.
Top global scores: \texttt{f10} (3.78), \texttt{f14} (1.76), \texttt{f18} (0.58), \texttt{f26} (0.57), \texttt{f4} (0.50).
\begin{itemize}\itemsep1pt
\item \texttt{f10}: type=concept, action=\textsc{candidate\_retire}.
\end{itemize}

\subparagraph{UCI\_wine.}
$n_{\mathrm{ref}}+n_{\mathrm{live}}=106+72$, shift columns $[4,\,8)$. Communication: 48 floats uplink (0.38 KiB) vs 2314 raw matrix floats (0.018 MiB if centralized)---\textbf{48$\times$} fewer scalars over the wire. FDR significant: 1; actions: \textsc{keep}=12, \textsc{candidate\\_retire}=1.
Top global scores: \texttt{f7} (2.86), \texttt{f1} (0.59), \texttt{f2} (0.54), \texttt{f0} (0.33), \texttt{f12} (0.04).
\begin{itemize}\itemsep1pt
\item \texttt{f7}: type=concept, action=\textsc{candidate\_retire}.
\end{itemize}

\subparagraph{UCI\_diabetes.}
$n_{\mathrm{ref}}+n_{\mathrm{live}}=265+177$, shift columns $[3,\,7)$. Communication: 39 floats uplink (0.30 KiB) vs 4420 raw matrix floats (0.034 MiB if centralized)---\textbf{113$\times$} fewer scalars over the wire. FDR significant: 1; actions: \textsc{keep}=9, \textsc{alert\\_recalibrate\\_and\\_review}=1.
Top global scores: \texttt{f6} (2.55), \texttt{f2} (0.66), \texttt{f0} (0.56), \texttt{f1} (0.25), \texttt{f3} (-0.67).
\begin{itemize}\itemsep1pt
\item \texttt{f6}: type=compound, action=\textsc{alert\_recalibrate\_and\_review}.
\end{itemize}

\subparagraph{UCI\_ionosphere.}
$n_{\mathrm{ref}}+n_{\mathrm{live}}=210+141$, shift columns $[12,\,20)$. Communication: 117 floats uplink (0.91 KiB) vs 11934 raw matrix floats (0.091 MiB if centralized)---\textbf{102$\times$} fewer scalars over the wire. FDR significant: 1; actions: \textsc{keep}=33, \textsc{alert\\_recalibrate\\_and\\_review}=1.
Top global scores: \texttt{f12} (4.87), \texttt{f5} (1.16), \texttt{f4} (0.84), \texttt{f2} (0.76), \texttt{f15} (0.62).
\begin{itemize}\itemsep1pt
\item \texttt{f12}: type=compound, action=\textsc{alert\_recalibrate\_and\_review}.
\end{itemize}

\subparagraph{UCI\_heart\_statlog.}
$n_{\mathrm{ref}}+n_{\mathrm{live}}=162+108$, shift columns $[4,\,9)$. Communication: 48 floats uplink (0.38 KiB) vs 3510 raw matrix floats (0.027 MiB if centralized)---\textbf{73$\times$} fewer scalars over the wire. FDR significant: 2; actions: \textsc{keep}=11, \textsc{alert\\_recalibrate\\_and\\_review}=2.
Top global scores: \texttt{f8} (2.68), \texttt{f5} (1.27), \texttt{f0} (0.57), \texttt{f3} (0.40), \texttt{f6} (0.19).
\begin{itemize}\itemsep1pt
\item \texttt{f8}: type=compound, action=\textsc{alert\_recalibrate\_and\_review}.
\item \texttt{f5}: type=compound, action=\textsc{alert\_recalibrate\_and\_review}.
\end{itemize}

\paragraph{Monitoring metrics (targets).} Drift-type accuracy under injection, FDR $\approx \alpha$, top-$k$ hit rate in shifted blocks, action precision, one-window closed-loop latency. Secure aggregation is orthogonal to Eff and left to production.
% === END federated_fsds_protocol_results.tex ===


\section{Communication-efficient cross-block monitoring}

\paragraph{Practical takeaway.}
Federated feature monitoring should treat \textbf{communication cost as a first-class constraint}.
Each client $k$ owns columns $F_k$ only; the server never receives raw rows $X$.
Instead, clients upload a \textbf{fixed-size summary per feature}---PO-risk distance proxy $d_f$, covariate-shift proxy $\Delta P_f$ (MMD-LOCO), and local rank---plus $O(1)$ block diagnostics (MMD$^2_k$, domain AUC$_k$, ESS$_k$).
Total uplink per round is $\sum_k O(|F_k|)$ scalars, versus $O(n\sum_k |F_k|)=O(np)$ if centralizing data.
This matches vertical federated learning while preserving the batch FSDS contract (REF vs LIVE, batch label $W$).

\paragraph{Communication envelope (per client $k$, one REF/LIVE window).}
\begin{center}
\small
\begin{tabular}{@{}lll@{}}
\toprule
\textbf{Message field} & \textbf{Size} & \textbf{Role} \\
\midrule
Block MMD$^2$, domain AUC, ESS & 3 floats & Gate / LOGO block score \\
Per $f\in F_k$: $(d_f, \Delta P_f, \text{rank}_k(f))$ & $3|F_k|$ floats & Global merge + FDR \\
Optional PO-risk $p$-value (block or global) & 0--1 float & Concept axis \\
\midrule
\textbf{Total uplink} & $3|F_k|+O(1)$ & No $n\times p$ matrix \\
\bottomrule
\end{tabular}
\end{center}
Let $B_k=3|F_k|+3$ be the float count for node $k$.
The \textbf{efficiency ratio} against naive centralization is
$\mathrm{Eff}= \frac{n(p+p)}{ \sum_k B_k } \approx \frac{2np}{\sum_k 3|F_k|}$ (same order as $2n/3$ when blocks partition columns evenly).

\paragraph{Privacy posture.}
Phase~2 supports plain upload (prototype) or a \texttt{secure\_aggregation} hook (masked sums / Shamir shares on \emph{aggregates only}).
Raw features never leave the client; only attribution statistics cross the wire.

\section{Explicit client--server protocol (three phases)}

\paragraph{Phase 1 --- Local attribution.}
On $F_k$ only, compute batch FSDS between REF and LIVE: domain RF VIMP $\Rightarrow d_f$ proxy; MMD-LOCO $\Rightarrow \Delta P_f$; rank within block.

\paragraph{Phase 2 --- Upload.}
Client sends \texttt{LocalAttributionPayload} (JSON or typed struct). Secure path: identical schema with masked values.

\paragraph{Phase 3 --- Global merge.}
Server (i) z-score or rank-combine $\{d_f\}$ across nodes; (ii) assign drift type per feature from $(d_f,\Delta P_f)$;
(iii) BH-FDR at level $\alpha$ on $p$-values derived from $d_f$; (iv) OFS actions:
\textsc{keep}, \textsc{recalibrate} (covariate + significant),
\textsc{candidate\_retire} (concept + significant),
\textsc{alert\_recalibrate\_and\_review} (compound + significant).

\section{Drift type and OFS/FDR closed loop}

\paragraph{Typing (feature level).}
High $\Delta P_f$, low $d_f$ $\Rightarrow$ covariate; high $d_f$, low $\Delta P_f$ $\Rightarrow$ concept; both high $\Rightarrow$ compound; else stable.

\paragraph{Closed loop per window.}
Drift diagnosis $\rightarrow$ $p$-values $\rightarrow$ FDR significant set $\rightarrow$ OFS actions $\rightarrow$ next window monitoring (online FDR extensions: SAFFRON/ADDIS on the merged $p$-stream).

\paragraph{Implementation.}
\texttt{Python/fsds\_sot/federated\_protocol.py}, \texttt{drift\_ofs\_fdr\_loop.py}; demo \texttt{demo\_federated\_protocol\_datasets.py}.

% === END federated_fsds_comprehensive_preamble.tex ===


\section{Federated protocol --- multi-dataset validation}
% === BEGIN federated_fsds_protocol_results.tex ===
% Auto-generated — demo_federated_protocol_datasets.py
\section{Empirical validation across tabular benchmarks}

\paragraph{Summary (communication + localization).}
Each run simulates vertical FL: REF/LIVE split, covariate injection on a contiguous feature block, then the three-phase protocol with BH-FDR ($\alpha{=}0.05$). \textbf{Eff} is the ratio of raw centralized matrix elements $(n_{\mathrm{ref}}+n_{\mathrm{live}})\times p$ to total uplink scalars $\sum_k(3|F_k|+3)$. \textbf{Top-1 hit} indicates whether the highest global attribution score falls inside the injected shift block.

\begin{center}
\small
\begin{tabular}{@{}lrrrrrrl@{}}
\toprule
Dataset & $p$ & nodes & FDR sig & Eff & uplink KiB & top-1 & hit \\
\midrule
WDBC\_breast\_cancer & 30 & 5 & 1 & 163$\times$ & 0.82 & \texttt{f10} & yes \\
UCI\_wine & 13 & 3 & 1 & 48$\times$ & 0.38 & \texttt{f7} & yes \\
UCI\_diabetes & 10 & 3 & 1 & 113$\times$ & 0.30 & \texttt{f6} & yes \\
UCI\_ionosphere & 34 & 5 & 1 & 102$\times$ & 0.91 & \texttt{f12} & yes \\
UCI\_heart\_statlog & 13 & 3 & 2 & 73$\times$ & 0.38 & \texttt{f8} & yes \\
\bottomrule
\end{tabular}
\end{center}

\subparagraph{WDBC\_breast\_cancer.}
$n_{\mathrm{ref}}+n_{\mathrm{live}}=341+228$, shift columns $[10,\,16)$. Communication: 105 floats uplink (0.82 KiB) vs 17070 raw matrix floats (0.130 MiB if centralized)---\textbf{163$\times$} fewer scalars over the wire. FDR significant: 1; actions: \textsc{keep}=29, \textsc{candidate\\_retire}=1.
Top global scores: \texttt{f10} (3.78), \texttt{f14} (1.76), \texttt{f18} (0.58), \texttt{f26} (0.57), \texttt{f4} (0.50).
\begin{itemize}\itemsep1pt
\item \texttt{f10}: type=concept, action=\textsc{candidate\_retire}.
\end{itemize}

\subparagraph{UCI\_wine.}
$n_{\mathrm{ref}}+n_{\mathrm{live}}=106+72$, shift columns $[4,\,8)$. Communication: 48 floats uplink (0.38 KiB) vs 2314 raw matrix floats (0.018 MiB if centralized)---\textbf{48$\times$} fewer scalars over the wire. FDR significant: 1; actions: \textsc{keep}=12, \textsc{candidate\\_retire}=1.
Top global scores: \texttt{f7} (2.86), \texttt{f1} (0.59), \texttt{f2} (0.54), \texttt{f0} (0.33), \texttt{f12} (0.04).
\begin{itemize}\itemsep1pt
\item \texttt{f7}: type=concept, action=\textsc{candidate\_retire}.
\end{itemize}

\subparagraph{UCI\_diabetes.}
$n_{\mathrm{ref}}+n_{\mathrm{live}}=265+177$, shift columns $[3,\,7)$. Communication: 39 floats uplink (0.30 KiB) vs 4420 raw matrix floats (0.034 MiB if centralized)---\textbf{113$\times$} fewer scalars over the wire. FDR significant: 1; actions: \textsc{keep}=9, \textsc{alert\\_recalibrate\\_and\\_review}=1.
Top global scores: \texttt{f6} (2.55), \texttt{f2} (0.66), \texttt{f0} (0.56), \texttt{f1} (0.25), \texttt{f3} (-0.67).
\begin{itemize}\itemsep1pt
\item \texttt{f6}: type=compound, action=\textsc{alert\_recalibrate\_and\_review}.
\end{itemize}

\subparagraph{UCI\_ionosphere.}
$n_{\mathrm{ref}}+n_{\mathrm{live}}=210+141$, shift columns $[12,\,20)$. Communication: 117 floats uplink (0.91 KiB) vs 11934 raw matrix floats (0.091 MiB if centralized)---\textbf{102$\times$} fewer scalars over the wire. FDR significant: 1; actions: \textsc{keep}=33, \textsc{alert\\_recalibrate\\_and\\_review}=1.
Top global scores: \texttt{f12} (4.87), \texttt{f5} (1.16), \texttt{f4} (0.84), \texttt{f2} (0.76), \texttt{f15} (0.62).
\begin{itemize}\itemsep1pt
\item \texttt{f12}: type=compound, action=\textsc{alert\_recalibrate\_and\_review}.
\end{itemize}

\subparagraph{UCI\_heart\_statlog.}
$n_{\mathrm{ref}}+n_{\mathrm{live}}=162+108$, shift columns $[4,\,9)$. Communication: 48 floats uplink (0.38 KiB) vs 3510 raw matrix floats (0.027 MiB if centralized)---\textbf{73$\times$} fewer scalars over the wire. FDR significant: 2; actions: \textsc{keep}=11, \textsc{alert\\_recalibrate\\_and\\_review}=2.
Top global scores: \texttt{f8} (2.68), \texttt{f5} (1.27), \texttt{f0} (0.57), \texttt{f3} (0.40), \texttt{f6} (0.19).
\begin{itemize}\itemsep1pt
\item \texttt{f8}: type=compound, action=\textsc{alert\_recalibrate\_and\_review}.
\item \texttt{f5}: type=compound, action=\textsc{alert\_recalibrate\_and\_review}.
\end{itemize}

\paragraph{Monitoring metrics (targets).} Drift-type accuracy under injection, FDR $\approx \alpha$, top-$k$ hit rate in shifted blocks, action precision, one-window closed-loop latency. Secure aggregation is orthogonal to Eff and left to production.
% === END federated_fsds_protocol_results.tex ===


\part{Reproducibility}

\section{Code and artifacts}
\begin{itemize}\itemsep3pt
  \item \textbf{Uplift two-batch:} \texttt{Python/loco\_auuc/}, \texttt{run\_uplift\_subset\_benchmark.py}, \texttt{plot\_uplift\_fsds\_monitor\_flow.py}.
  \item \textbf{Artifacts:} \texttt{artifacts/uplift\_two\_layer\_benchmark.json}, \texttt{artifacts/uplift\_fsds\_monitor\_flow\_en.png}.
  \item \textbf{Federated protocol:} \texttt{Python/fsds\_sot/federated\_protocol.py}, \texttt{drift\_ofs\_fdr\_loop.py}, \texttt{demo\_federated\_protocol\_datasets.py}.
  \item \textbf{Artifacts:} \texttt{artifacts/federated\_protocol\_benchmark.json}.
  \item \textbf{Playbooks:} \texttt{docs/FSDS\_UPLIFT\_TWO\_BATCH\_LANDING.md}, \texttt{docs/FSDS\_FEDERATED\_FEATURE\_BLOCKS.md}.
\end{itemize}

\section{Included \texttt{.tex} fragments (maintainers)}
\begin{itemize}\itemsep2pt
  \item \texttt{uplift\_fsds\_two\_batch\_formulation.tex}
  \item \texttt{uplift\_fsds\_benchmark\_results.tex} (auto)
  \item \texttt{uplift\_fsds\_subset\_localization\_results.tex} (auto)
  \item \texttt{uplift\_fsds\_subset\_insights.tex}
  \item \texttt{federated\_fsds\_comprehensive\_preamble.tex}
  \item \texttt{federated\_fsds\_protocol\_results.tex} (auto)
\end{itemize}

\end{document}
